%%%%%%%%%%%%%%%%%%%%%%%%%%%%%%%%%
%
%       EXERCÍCIO
%
%%%%%%%%%%%%%%%%%%%%%%%%%%%%%%%%%

\ifdefstring{\atividade}{prova}{%
    \ifdefstring{\modo}{objetivo}{%
        \renewcommand{\valorquestao}{\ValorQObj\ ponto}
    }{%
        \renewcommand{\valorquestao}{\ValorQDisc\ pontos}
    }%
}%

\begin{exercicioBanco}[\valorquestao]
Sejam \(A = \{-2,-1,0,1,2\}\) e \(B = \{x \in \bbZ \mid -5 \le x \le 5\}\). Considere a função
%
\[
    \begin{array}{rcl}
        f: \; A & \to & B \\
        x & \mapsto & x^{2} + 1.
    \end{array}
\]
%
Determine o conjunto imagem de \(f\).

% Define as alternativas
\newcommand{\alternativas}{%
    \begin{center}
        \begin{tabularx}{\textwidth}{XXX}
            (a) \(\{1,2,5\}\). &
            (b) \(\{-5,-2,1,2,5\}\). &
            (c) \(A\). \\[5pt]
            (d) \(\{0,1,4\}\). &
            (e) \(\{-1,0,3\}\).
        \end{tabularx}
    \end{center}
}

% Define a resposta correta
\newcommand{\resposta}{A}

% Lógica condicional para exibição
\ifdefstring{\atividade}{lista}{%
    \alternativas
    \vspace{0.5em}
    
    \noindent\textbf{Resposta:} letra \textbf{\resposta}.
}{%
    \ifdefstring{\modo}{objetivo}{%
        \alternativas
    }{%
        % modo = discursiva → não mostra alternativas
    }
}
\end{exercicioBanco}

